% Generated by scripts/assemble.py; corrections belong in source/corrections.json.
\examhead{95}{河海大学 2026 年数学分析试卷}{94}
% p095-b09; source PDF page 95
\par\Needspace{5\baselineskip}
\blockanchor{p095-b09}
\majorquestion{一}判断题. 每题 5 分, 共 20 分.\par
% p095-b10; source PDF page 95
\blockanchor{p095-b10}
\question{1}设 $\lambda > 1$ 及 $\{a_n\}$ 是实数列, 如果对所有 $n, p \in \mathbb{N}^+$ 都有 $|a_{n+p} - a_n| < \frac{p}{n^\lambda}$, 那么 $\{a_n\}$ 是基本数列(Cauchy 列).\par
% p095-b11; source PDF page 95
\blockanchor{p095-b11}
\question{2}设 $x_0$ 是连续函数 $f(x)$ 的极小值点, 则存在 $\delta > 0$, 使得 $f(x)$ 在 $(x_0 - \delta, x_0]$ 上单调递减, 在 $[x_0, x_0 + \delta)$ 上单调递增.\par
% p095-b12; source PDF page 95
\blockanchor{p095-b12}
\question{3}设 $f(x)$ 在 $[a, b]$ 上黎曼可积, 则存在可导函数 $F(x)$, 使得 $F'(x) = f(x)$, $x \in [a, b]$.\par
% p095-b13; source PDF page 95
\blockanchor{p095-b13}
\question{4}设实数列  $\{a_{n}\}$  满足  $a_{n}>0\ (n=1,2,\cdots)$ ，且  $\lim_{n\to\infty}a_{n}=0$ ，则级数  $\sum_{n=1}^{\infty}(-1)^{n-1}a_{n}$  收敛.\par
% p095-b14; source PDF page 95
\par\Needspace{5\baselineskip}
\blockanchor{p095-b14}
\majorquestion{二}叙述题. 每题 5 分, 共 20 分.\par
% p095-b15; source PDF page 95
\blockanchor{p095-b15}
\question{5}写出“当 $x \to +\infty$ 时，函数 $f(x)$ 的极限为负无穷大”的否定命题的分析表述.\par
% p095-b16; source PDF page 95
\blockanchor{p095-b16}
\question{6}叙述判断反常积分 $\int_{a}^{+\infty} f(x) g(x) \, \mathrm{d}x$ 收敛的 Abel 判别法.\par
% p095-b17; source PDF page 95
\blockanchor{p095-b17}
\question{7}叙述判定函数项级数 $\sum_{n=1}^{\infty} u_n(x)$ 关于 $x$ 在集合 $D$ 上一致收敛的 Cauchy 收敛原理.\par
% p095-b18; source PDF page 95
\blockanchor{p095-b18}
\question{8}叙述二元函数 $z = z(x, y)$ 在点 $(x_0, y_0)$ 处可微的定义.\par
% p095-b19; source PDF page 95
\par\Needspace{5\baselineskip}
\blockanchor{p095-b19}
\majorquestion{三}解答与证明. 每题 10 分, 共 50 分.\par
% p095-b20; source PDF page 95
\blockanchor{p095-b20}
\question{9}计算极限 $\lim_{x\to +\infty}\left(\sqrt[3]{x^3 + 2x^2 + 1} -\sqrt[5]{x^5 + 3x^4 - 2}\right).$\par
% p095-b21; source PDF page 95
\blockanchor{p095-b21}
\question{10}计算不定积分 $\int \frac{\mathrm{d}x}{\sqrt{\sin^7 x \cos x}}$.\par
% p095-b22; source PDF page 95
\blockanchor{p095-b22}
\question{11}求幂级数 $\sum_{n=0}^{\infty} \frac{x^n}{4n^2 - 1}$ 的收敛域及和函数.\par
% p095-b23; source PDF page 95
\par\Needspace{6\baselineskip}
\blockanchor{p095-b23}
\question{12}设 $\alpha > 0$, 且 $f(x, y)$ 是 $\mathbb{R}^2$ 上的可微函数. 证明: 对任意 $t > 0$ 及 $(x, y) \in \mathbb{R}^2 \setminus \{(0, 0)\}$, 成立 $f(tx, ty) = t^\alpha f(x, y)$ 的充要条件是对任意的 $(x, y) \in \mathbb{R}^2$, 成立\par
\nopagebreak[4]
% p095-b24; source PDF page 95
\blockanchor{p095-b24}
\displayformula{x \frac {\partial f}{\partial x} (x, y) + y \frac {\partial f}{\partial y} (x, y) = \alpha f (x, y).}
% p096-b00; source PDF page 96
\par\Needspace{6\baselineskip}
\blockanchor{p096-b00}
\question{13}计算曲面积分\par
\nopagebreak[4]
% p096-b01; source PDF page 96
\blockanchor{p096-b01}
\displayformula{\iint_ {\Sigma} \frac {x \mathrm{d} y \mathrm{d} z + y \mathrm{d} z \mathrm{d} x + z \mathrm{d} x \mathrm{d} y}{(x ^ {2} + y ^ {2} + z ^ {2}) ^ {\frac {3}{2}}}.}
% p096-b02; source PDF page 96
\blockanchor{p096-b02}
其中 $\Sigma$ 是曲面 $9 - z = \sqrt{(x - 1)^2 + (y - 2)^2} (z \geq 0)$, 方向取上侧.\par
% p096-b03; source PDF page 96
\par\Needspace{5\baselineskip}
\blockanchor{p096-b03}
\majorquestion{四}证明题. 每题 15 分, 共 60 分.\par
% p096-b04; source PDF page 96
\blockanchor{p096-b04}
\question{14}设 $f(x)$ 和 $g(x)$ 在 $[a, b]$ 上二阶可导，且 $g''(x) \neq 0$，$f(a) = f(b) = g(a) = g(b) = 0$.\par
% p096-b05; source PDF page 96
\blockanchor{p096-b05}
(1) 证明: 对任意的 $x \in (a, b)$, 都有 $g(x) \neq 0$.\par
% p096-b06; source PDF page 96
\blockanchor{p096-b06}
(2) 证明: 存在 $\xi \in (a, b)$, 使得 $\frac{f(\xi)}{g(\xi)} = \frac{f''(\xi)}{g''(\xi)}$.\par
% p096-b07; source PDF page 96
\blockanchor{p096-b07}
\question{15}设 $\alpha > 0$.\par
% p096-b08; source PDF page 96
\blockanchor{p096-b08}
(1) 证明: 反常积分 $\int_{1}^{+\infty} \frac{\mathrm{d}x}{(1 + x^2 y^2) x^\alpha}$ 关于 $y$ 在 $(0, +\infty)$ 上内闭一致收敛.\par
% p096-b09; source PDF page 96
\par\Needspace{6\baselineskip}
\blockanchor{p096-b09}
(2) 证明: $I(y) = \int_{1}^{+\infty} \frac{\arctan(xy)}{x^{\alpha + 1}} \, \mathrm{d}x$ 是 $(0, +\infty)$ 上的可导函数, 且满足方程\par
\nopagebreak[4]
% p096-b10; source PDF page 96
\blockanchor{p096-b10}
\displayformula{I ^ {\prime} (y) - \alpha I (y) + \arctan y = 0.}
% p096-b11; source PDF page 96
\par\Needspace{6\baselineskip}
\blockanchor{p096-b11}
\question{16}设 $f(x)$ 在 $[a, b]$ 上连续，且满足 $0 < m \leq f(x) \leq M$，证明：\par
\nopagebreak[4]
% p096-b12; source PDF page 96
\blockanchor{p096-b12}
\displayformula{1 \leq \frac {1}{(b - a) ^ {2}} \int_ {a} ^ {b} f (x) \mathrm{d} x \int_ {a} ^ {b} \frac {\mathrm{d} x}{f (x)} \leq \frac {(m + M) ^ {2}}{4 m M}.}
% p096-b13; source PDF page 96
\par\Needspace{6\baselineskip}
\blockanchor{p096-b13}
\question{17}设 $a_{n} \neq 0 (n = 1,2,\dots)$, 对级数 $\sum_{n=1}^{\infty} a_{n}$ 加括号得到\par
\nopagebreak[4]
% p096-b14; source PDF page 96
\blockanchor{p096-b14}
\displayformula{\left(a _ {1} + a _ {2} + \dots + a _ {n _ {1}}\right) + \left(a _ {n _ {1} + 1} + a _ {n _ {1} + 2} + \dots + a _ {n _ {2}}\right) + \dots + \left(a _ {n _ {k - 1} + 1} + a _ {n _ {k - 1} + 2} + \dots + a _ {n _ {k}}\right) + \dots .}
% p096-b15; source PDF page 96
\par\Needspace{6\baselineskip}
\blockanchor{p096-b15}
其中每个括号内的 $a_{n}$ 具有相同的符号，相邻括号内 $a_{n}$ 的符号不同. 记 $n_0 = 0$ ，且\par
\nopagebreak[4]
% p096-b16; source PDF page 96
\blockanchor{p096-b16}
\displayformula{b _ {k} = \sum_ {j = n _ {k - 1} + 1} ^ {n _ {k}} a _ {j} (k = 1, 2, \dots).}
% p096-b17; source PDF page 96
\blockanchor{p096-b17}
(1) 证明: $\sum_{n=1}^{\infty} a_n$ 和 $\sum_{n=1}^{\infty} b_n$ 具有相同的敛散性.\par
% p096-b18; source PDF page 96
\blockanchor{p096-b18}
(2) 证明: $\sum_{n=1}^{\infty} \frac{(-1)^{[\sqrt{n}]}}{n}$ 是收敛的, 其中 $[x]$ 是取整函数, 表示不超过 $x$ 的最大整数.\par
